\documentclass[../../main.tex]{subfiles}
\begin{document}

\section{切空间的其他定义}\label{section:ch03-切空间的其他定义}

在文献中，你会发现光滑流形的切向量以几种不同方式定义。这里我们描述最常见的几种。（习题~11-4 后的注记提示了另一种定义。）熟悉所有这些是有益的。本节中，$M$ 表示任意一个带边或无边光滑流形。

\subsection{作为芽空间导子的切向量}\label{subsection:ch03-作为芽导子的切向量}

最常见的另一种定义基于光滑函数"芽"的概念，我们现在定义它。

$M$ 上的一个光滑函数元素是一个有序对 $(f,U)$，其中 $U$ 是 $M$ 的一个开子集，$f:U\to\mathbb{R}$ 是一个光滑函数。给定一点 $p\in M$，让我们在所有定义域包含 $p$ 的光滑函数元素所成集合上定义一个等价关系，规定 $(f,U)\sim(g,V)$，若 $f\equiv g$ 在 $p$ 的某个邻域上。一个函数元素 $(f,U)$ 的等价类称为 $f$ 在 $p$ 处的芽。$p$ 处所有光滑函数芽所成集合记作 $C^{\infty}_{p}(M)$。它在运算
\begin{align*}
c\cdot[(f,U)]=[(cf,U)],\quad [(f,U)]+[(g,V)]=[(f+g,U\cap V)],\quad [(f,U)]\cdot[(g,V)]=[(fg,U\cap V)]
\end{align*}
下是一个实向量空间和一个结合代数。（这个代数的零元是 $M$ 上零函数的等价类。）

让我们把函数元素 $(f,U)$ 在 $p$ 处的芽简记为 $[f]_{p}$；不必在记号中包含定义域 $U$，因为同一个芽由 $f$ 在 $p$ 任意邻域上的限制所代表。说两个芽 $[f]_{p}$ 和 $[g]_{p}$ 相等，只不过是说 $f\equiv g$ 在 $p$ 的某个无论多小的邻域上。

$C^{\infty}_{p}(M)$ 的一个导子是一个线性映射 $v:C^{\infty}_{p}(M)\to\mathbb{R}$，满足与~\eqref{align:ch03-流形导子} 类似的如下乘积法则：
\begin{align*}
v[fg]_{p}=f(p)v[g]_{p}+g(p)v[f]_{p}.
\end{align*}
通常把 $M$ 在 $p$ 处的切空间定义为 $C^{\infty}_{p}(M)$ 的导子所成向量空间 $D_{p}M$。多亏命题~\cref{proposition:ch03-切向量局部作用}，证明 $D_{p}M$ 与我们所定义的切空间自然同构是一件简单的事（见习题~3-7）。

芽的定义有一些优点。其中最重要的之一是，它使切空间的局部性更清楚，而无需使用鼓包函数。因为不存在解析鼓包函数，切向量的芽定义是实解析或复解析流形上唯一可用的定义。芽方法的主要缺点只不过是，它给一个已经高度抽象的定义又增加了一层复杂性。

\subsection{作为曲线等价类的切向量}\label{subsection:ch03-作为曲线等价类的切向量}

另一种常见的切向量方法，是在具有相同起点的光滑曲线集合上定义一个内在等价关系，它抓住"具有相同速度"这一想法，并把切向量定义为曲线的等价类。这里我们描述这样一个等价关系。

设 $p$ 是 $M$ 的一点。我们希望在所有形如 $\gamma:J\to M$（其中 $J$ 是一个包含 $0$ 的区间、$\gamma(0)=p$）的光滑曲线集合上定义一个等价关系。给定两条这样的曲线 $\gamma_{1}:J_{1}\to M$ 和 $\gamma_{2}:J_{2}\to M$，让我们说 $\gamma_{1}\sim\gamma_{2}$，若对每个定义在 $p$ 邻域上的光滑实值函数 $f$ 都有 $(f\circ\gamma_{1})'(0)=(f\circ\gamma_{2})'(0)$。令 $V_{p}M$ 记这些等价类所成集合。$M$ 在 $p$ 处的切空间常被定义为集合 $V_{p}M$。

用这个定义，定义一个光滑映射 $F:M\to N$ 的微分非常容易：它是那个把 $[\gamma]\in V_{p}M$ 送到 $[F\circ\gamma]\in V_{F(p)}N$ 的映射。光滑曲线的速度向量也几乎同样容易定义。设 $\gamma:J\to M$ 是任意光滑曲线。若 $0\in J$，则 $\gamma$ 在 $0$ 处的速度正是 $\gamma$ 在 $V_{\gamma(0)}M$ 中的等价类。任意其他点 $t_{0}\in J$ 处的速度，可以定义为由 $\gamma_{t_{0}}(t)=\gamma(t_{0}+t)$ 定义的曲线 $\gamma_{t_{0}}$ 在 $V_{\gamma(t_{0})}M$ 中的等价类。

习题~3-8 表明，在 $V_{p}M$ 与 $T_{p}M$ 之间有一个自然的一一对应。这个定义的优点是几何上更直观，但它有一个严重缺点：$V_{p}M$ 上向量空间结构的存在一点也不显然。

\subsection{作为 $n$ 元组等价类的切向量}\label{subsection:ch03-作为n元组等价类的切向量}

定义切空间的又一种方法，基于切向量分量在坐标中的变换规则~\eqref{align:ch03-分量变换}。人们把 $p\in M$ 处的一个切向量定义为这样一条规则：它给包含 $p$ 的每个光滑坐标图指派一个有序 $n$ 元组 $(v^{1},\dots,v^{n})\in\mathbb{R}^{n}$，并具有指派给重叠坐标图的 $n$ 元组按~\eqref{align:ch03-分量变换} 变换的性质。（事实上，这是所有定义中最古老的一个，许多物理学家仍倾向于这样思考切向量。）

在这种方法中，一条曲线的速度由坐标中通常的欧几里得公式定义，而 $F:M\to N$ 的微分被定义为由 $F$ 在坐标中的 Jacobian 矩阵确定的线性映射。然后人们必须借助涉及链式法则的冗长计算，证明这些运算良定义、独立于坐标的选取。

选择 $T_{p}M$ 的各种刻画中哪一个作为定义，是个人品味的问题。然而，我们所选择的定义，无论起初看起来多么抽象，都有几个优点：它相对具体（切向量是 $C^{\infty}(M)$ 的真实导子，不涉及等价类）；它使 $T_{p}M$ 上的向量空间结构显而易见；它导向微分、速度以及我们将研究的许多其他几何对象的直接、与坐标无关的定义。

\end{document}